\section{Data and prediction task}
\label{sec:data}
We use the MIDI-only MAESTRO v3.0.0 piano performances \citep{hawthorne2019maestro}. The supplied training/validation/test splits contain 962/137/177 pieces and 33,293,600/3,790,187/4,334,269 event tokens. The documentation keeps performances of the same composition in one split \citep{maestrodata}; we preserve this split but do not independently audit every possible overlap.

The tokenizer converts music into 376 instructions: 128 note starts, 128 note endings, 100 time shifts, 16 velocity bins, two pedal states, and BOS/EOS boundaries. Timing is rounded to 10\,ms; longer pauses use multiple shifts of at most one second. At equal timestamps, note endings precede pedal changes and note starts. A velocity token precedes every note start. Only pedal state changes are emitted. This preserves expressive event order while discarding some timing/velocity detail and other MIDI metadata. A token is not a note, and 1,024 tokens cover different musical durations in busy and quiet passages.

The notebook checks token ranges, lengths, and piece boundaries and decodes one source piece to inspect the tokenizer. This is not a generated sample. Storage details and preparation steps are in the reproduction guide; our results concern event prediction, not audio reconstruction.
